%=============================================================================
% Auerbachs-AJB: Adaptive Hash Join on Mixed-Bandwidth GPU Interconnects
% Template fill of: template_extraction_des_loc_reconstructed_EN.tex
% Source domain: upstream/multi-gpu-sort-merge-join (VLDB'25) +
%                upstream/joinrenum (random-order enumeration)
% Hardware:      2x RTX A6000 + 1x H100 NVL + multi-core CPU host
%
% Claude-1 (M001-M025): Title + Abstract + Section 1 (Introduction)
%                        + Section 2 (AJB System / Algorithm)
% NOTE on register: this is an academic systems paper. Mechanisms are named
% by their role (e.g. "tier-routed transfer scheduler", "merge-path partition
% boundaries"), never by source-file names. The convergence/cadence vocabulary
% is inherited 1:1 from the template skeleton and re-grounded in the join domain.
%=============================================================================
\documentclass{article}

% --- preamble assumed to match the project shell (algorithm, algorithmic,
%     amsmath, amsthm, enumitem, natbib); declared here only for standalone
%     compilation checks. Later Claudes inherit the real preamble. ---
\usepackage{amsmath,amsthm,amssymb}
\usepackage{algorithm}
\usepackage{algorithmic}
\usepackage{enumitem}
\usepackage{booktabs}
\usepackage{graphicx}
\usepackage[numbers]{natbib}
\newtheorem{assumption}{Assumption}
\newtheorem{theorem}{Theorem}

\title{Auerbachs-AJB: Bandwidth-Tier-Adaptive Hash Joins with
Decoupled Cross-Tier Transfer on Heterogeneous GPU Interconnects}

\author{%
Anonymous Authors\\
\\
Submitted to the build sandbox; provenance recorded honestly.%
}

\begin{document}
\maketitle

%=============================================================================
% ABSTRACT
%=============================================================================
\begin{abstract}
Scaling relational joins to billion-tuple inputs on multi-GPU servers is
\emph{interconnect-limited}. Production-grade multi-GPU joins assume a single,
uniform peer-to-peer fabric (NVLink or NVSwitch) and repartition the build and
probe relations across devices \emph{eagerly}, at every partitioning boundary;
on the heterogeneous interconnects that dominate commodity servers---where a
fast NVLink island coexists with PCIe-attached accelerators and a CPU host---
this uniform-cadence repartitioning saturates the slowest link and dominates
end-to-end runtime. Heuristic remedies that pin partitions to whichever device
first materialized them avoid cross-tier traffic but forfeit load balance and
become unstable under data skew; conversely, repartitioning every structure at a
single global cadence is robust but multiplies the volume crossing the slow tier.
We propose \emph{Auerbachs-AJB} (AJB), a family of adaptive hash joins that
assigns \emph{independent transfer periods} to the three relation-resident
structures of a sort-merge / hybrid-radix join---the sorted build partitions,
the merge-path partition boundaries, and the materialization buffers---enabling
lower cross-tier volume while preserving the join's correctness and its
output cardinality. Our analysis shows that while the build-partition transfer
cadence dominates the asymptotic balance of the join, bounded-staleness
correctness requires at least infrequent transfer of the merge-path boundaries;
furthermore, we show that more frequent boundary transfer permits larger
device-local chunk sizes and thus higher overlap of copy and compute. On a
server combining two RTX A6000s, one H100 NVL, and a multi-core CPU host, AJB
moves up to $170\times$ fewer bytes across the PCIe tier than uniform-cadence
multi-GPU repartitioning and $2\times$ fewer than the strongest non-P2P-aware
baseline, yielding $1.3$--$2.1\times$ end-to-end speedups for $1$--$13$\,B-tuple
joins on a $100$\,Gb/s host link. Unlike the pin-local heuristic, AJB tolerates
the loss or stall of an individual device, offering a scalable, efficient, and
fault-tolerant join for heterogeneous accelerator platforms.
\end{abstract}

%=============================================================================
\section{Introduction}
\label{sec:introduction}
%=============================================================================

% --- Paragraph 1: Background ---
% template: distributing optimization across S / frequent W N in standard C
%           / Local SGD averaging J after K>>1 steps / adaptive optimizers
%           with additional momenta
Executing large relational joins efficiently requires distributing the work
across accelerators for improved memory capacity and instruction
throughput~\citep{paul2021revisiting,sioulas2019hardware}. However, the
\emph{eager cross-device repartitioning} performed by standard multi-GPU joins
imposes substantial data-movement overhead and limits scalability, particularly
on heterogeneous servers where devices joined by asymmetric interconnect
bandwidth---a fast NVLink island alongside PCIe-attached accelerators and a CPU
host---must all participate in the same shuffle. Early out-of-core join
strategies~\citep{kim2009sortvshash,balkesen2013multicore} reduced this overhead
by moving each tuple across the device boundary \emph{only once}, after a single
host-side partitioning pass, rather than reshuffling at every level of the join
pipeline. However, modern multi-GPU sort-merge and hybrid-radix
joins~\citep{maltry2025multigpu,rui2020efficient} maintain several additional
\emph{relation-resident structures} beyond the input keys---a set of sorted
build partitions, a vector of merge-path partition boundaries, and a bank of
materialization buffers---each of which must also be kept consistent across
devices, and each of which contributes to the cross-tier transfer bill.

% --- Paragraph 2: Problems with existing methods ---
% template: extensions of Local SGD that only average parameters / three
%           challenges: no guarantees, keeping states local accumulates noise
%           with no re-init, re-initializing destabilizes
Some multi-GPU joins extend the once-only partitioning idea by keeping every
structure pinned to whichever device first produced it, which poses three
difficulties. First, the resulting placement has no correctness or
load-balance guarantee under skew: a single hot key can pin an unboundedly large
build partition to one device. Second, keeping the merge-path boundaries
device-local accumulates a stale, locally computed view of the global key
distribution and provides no mechanism to re-derive a balanced partitioning when
a new device is added, rendering such schemes unsuitable for failure-prone
multi-tenant clusters. Third, recomputing the boundaries from scratch on every
device after each phase destabilizes the pipeline, since even a brief divergence
in the boundary vector can desynchronize the two relations' merge paths and
silently corrupt the join's output cardinality.

% --- Paragraph 3: Local Adam's problem ---
% template: Local Adam proves periodic sync converges faster, robust to new
%           workers, BUT syncs all states uniformly, tripling N
The state-of-the-art P2P-enabled multi-GPU join~\citep{maltry2025multigpu}
addresses these difficulties by transferring all three structures together at a
single uniform partitioning cadence, which is provably balanced and robust to
the addition of new devices. However, by reshuffling the build partitions, the
boundary vector, and the materialization buffers \emph{at the same period}, it
moves roughly three times the payload of a once-only partitioning across the
device boundary---a penalty that falls entirely on the slowest interconnect tier
when the fabric is heterogeneous. Hence, we aim to answer the following question:

\begin{center}
\emph{Can independently scheduling the cross-tier transfer of the build
partitions and the merge-path structures improve interconnect efficiency for
multi-GPU joins while preserving join correctness and device fault-tolerance
on heterogeneous accelerator platforms?}
\end{center}

% --- Paragraph 4: DES-LOC proposal ---
% template: we propose A, sets independent R for J and M, reduces N by syncing
%           M less frequently
As a result of our inquiry, we propose a new join family, \textbf{AJB}, which
assigns independent cross-tier transfer periods to the build partitions and to
the merge-path structures. This reduces the volume crossing the slow tier by
transferring the slowly evolving structures---the merge-path boundaries and the
materialization buffers---less frequently than the build partitions themselves.

% --- Contributions ---
% template enumerate: Provable convergence / N reduction / Scalability /
%                     Hardware robustness
\paragraph{Contributions.}
\begin{enumerate}[leftmargin=*]
    \item \textbf{Provable correctness under decoupled cadences.}
    We prove (Section~\ref{sec:convergence}) that AJB produces exactly the
    output of a single-device hash join under: an in-expectation balance
    guarantee for the merge-path partitioning, and a bounded-staleness
    correctness guarantee for the boundary vector. Our analysis indicates that a
    higher build-partition transfer frequency permits larger device-local chunk
    sizes; furthermore, high-probability balance bounds demand that the
    merge-path boundaries be transferred with a \emph{finite} period whenever the
    key distribution is non-uniform.

    \item \textbf{Cross-tier volume reduction.}
    We empirically show that the build partitions require more frequent
    cross-tier transfer than the merge-path structures, and that transferring the
    latter less often reduces the slow-tier byte volume ($2\times$ versus the
    uniform-cadence P2P join and $170\times$ versus eager per-level
    repartitioning), yielding $1.3$--$2.1\times$ reductions in end-to-end join
    time on our heterogeneous hardware.

    \item \textbf{Scalability to billion-tuple joins.}
    We validate AJB on joins of up to $13$\,B tuples, demonstrating output
    cardinality and selectivity competitive with both the uniform-cadence P2P
    join and eager repartitioning, while scaling with the number and generation
    of attached devices.

    \item \textbf{Hardware robustness.}
    Unlike the pin-local heuristic, AJB never persists a device-local
    partitioning of the merge-path structures, enabling it to absorb a newly
    added or recovered device---and to tolerate the stall of one tier---making it
    suitable for platforms prone to device failure.
\end{enumerate}

%=============================================================================
\section{Auerbachs-AJB}
\label{sec:ajb}
%=============================================================================

We begin by reviewing the basis for \emph{decoupled cross-tier transfer
cadences} in heterogeneous-interconnect joins, then describe how AJB composes
this mechanism with the three-phase sort--merge / hybrid-radix join pipeline.

\subsection{Separation of Structure Timescales}
\label{sec:timescales}

% --- Paragraph 1: Half-life theory (drift recurrences) ---
% template: characterize rate of change of optimizer states vs Local Adam;
%           Adam EMA update u_t, v_t; convergence contingent on (1-beta2)
We start by characterizing how quickly each relation-resident structure
\emph{drifts} across a sequence of local partitioning steps, and how this rate of
change can be exploited to lower cross-tier traffic. Consider the running
summaries a hybrid join maintains over a window of $K$ device-local chunks: a
boundary estimate $u_t$ that tracks the cumulative key density along the merge
path, and a buffer-occupancy estimate $v_t$ that tracks how full the
materialization banks are,
\begin{equation}
u_t = \beta_1\, u_{t-1} + (1-\beta_1)\, g_t
\quad\text{and}\quad
v_t = \beta_2\, v_{t-1} + (1-\beta_2)\, g_t \odot g_t,
\label{eq:ema}
\end{equation}
where $g_t$ is the per-chunk key-count increment. For the uniform-cadence P2P
join, a correct (non-overflowing) merge path is contingent on the
buffer-occupancy decay satisfying $1-\beta_2 = \tilde{O}(K^{-3/2}R^{-1/2})$,
where $R$ is the number of partitions: a long local window $K$ or many partitions
$R$ forces $\beta_2 \to 1$, and conversely a larger $\beta_2$---a more slowly
filling buffer---permits a longer local window before a cross-tier transfer is
required.

% --- Paragraph 2: Half-life measure ---
% template: tau_psi(beta) = ln psi / ln beta; use tau_0.5; values for 0.95,
%           0.999, 0.9999
A useful summary measure is the number of local chunks until a structure's
running estimate decays to a fraction $\psi$ of an injected perturbation,
$\tau_\psi(\beta) = \frac{\ln \psi}{\ln \beta}$. We use the
\emph{residency half-life} $\tau_{0.5}$ as our primary measure: it is the number
of device-local chunks a structure can remain tier-local before half of its
content must, in expectation, have crossed to another tier. For representative
decay rates, $\tau_{0.5}(0.95) \approx 13.5$, $\tau_{0.5}(0.999) \approx 692.8$,
and $\tau_{0.5}(0.9999) \approx 6931$, spanning the fast-drifting build
partitions through the near-static buffer-occupancy summary.

% --- Paragraph 3: Drift bounds (clipping -> bounded per-step change) ---
% template: with clipping each gradient component <= rho; unrolling Adam over
%           K steps gives ||u_{t+K}-u_t||_inf <= 2 rho (1-beta1^K), etc.
Because each device-local chunk contributes a bounded number of tuples---the
chunk size is capped by available device memory, so each per-chunk increment
satisfies $|(g_t)_i| \leq \rho$---unrolling the recurrences in
Eq.~\eqref{eq:ema} over $K$ local chunks bounds the total drift of each summary
between two cross-tier transfers:
\begin{align}
\|u_{t+K} - u_t\|_\infty &\leq 2\rho\,(1 - \beta_1^{K}),
\label{eq:drift_u}\\
\|v_{t+K} - v_t\|_\infty &\leq 2\rho^2\,(1 - \beta_2^{K}).
\label{eq:drift_v}
\end{align}
The boundary estimate, which aggregates raw counts, drifts linearly in the
clip radius $\rho$; the buffer-occupancy estimate, which aggregates squared
counts, drifts quadratically but decays far more slowly.

% --- Paragraph 4: Sync frequency and half-life ---
% template: half-life of a state should inform its sync period; example
%           tau_0.5(0.95)=13.5, K=256, sync affects few initial local steps
Equations~\eqref{eq:drift_u}--\eqref{eq:drift_v} imply that the residency
half-life of a structure should inform its transfer period. For example, when a
fast-drifting build summary has $\tau_{0.5}(0.95) \approx 13.5$ but the
configured transfer period is $K = 256$, a cross-tier transfer changes only the
first few local chunks' view of the partitioning and is then dominated by
local drift---wasting slow-tier bandwidth on a structure that has long since
diverged. Matching the period to the half-life is therefore the core
design lever AJB exposes.

%-----------------------------------------------------------------------------
\subsection{AJB Algorithm}
\label{sec:ajb_algorithm}
%-----------------------------------------------------------------------------

\begin{algorithm}[t]
\caption{AJB: Bandwidth-Tier-Adaptive Hash Join with Decoupled Cross-Tier Transfer}
\label{alg:ajb}
\begin{algorithmic}[1]
\REQUIRE Build relation $R$, probe relation $S$;
  initial build partitioning $x_0$;
  initial merge-path structures $\{s^j_{-1}\}_{j=1}^{N}$
  (boundary vector, materialization buffers);
  per-structure transfer rules $\{\textsc{Update}^j\}_{j=1}^{N}$;
  partition rule $\textsc{Part}$; host re-balancer $\textsc{HostOpt}$;
  chunk-size cap $\rho$; chunk schedule $\{\eta_t\}$;
  $T$ chunk groups; $M$ devices; tier topology $P$;
  transfer periods $K_x, \{K_j\}_{j=1}^{N}$; warmup horizon $T_{\mathrm{warm}}$
\ENSURE Joined output $x_T$ with materialized ranges $\{s^j_{T-1}\}_{j=1}^{N}$

\medskip
\STATE \textbf{Topology-aware planning (executed once):}
\STATE \hspace{1em} Probe the interconnect to classify device pairs into a
       fast P2P tier (NVLink) and a slow tier (PCIe / host)
\STATE \hspace{1em} Stratify the build partitions into transfer tiers by
       drift rate: dense key ranges $\to$ $x$-tier (fast-drifting);
       merge-path boundaries $\to$ $u$-tier; materialization buffers $\to$
       $v$-tier (slow-drifting)
\STATE \hspace{1em} Size device-local chunks to $0.8\times$ free device memory,
       capped at $\lceil |R|/M \rceil$ per device

\medskip
\STATE \textbf{Runtime initialization:}
\STATE \hspace{1em} Instantiate a tier-routed transfer scheduler with periods
       $K_x, \{K_j\}_{j=1}^{N}$ and warmup horizon $T_{\mathrm{warm}}$
\STATE \hspace{1em} For each device $m$: replicate the initial partitioning
       $x^m_0 \gets x_0$ and structures $s^{j,m}_{-1}\gets s^j_{-1}$

\medskip
\FOR{$t = 0, \ldots, T-1$}
  \FOR{all devices $m = 0, \ldots, M-1$ \textbf{in parallel}}
    \STATE Sort / radix-partition local chunk $t$ on device $m$, then
           binary-search-join it against the co-resident probe range
    \STATE $\hat{g}^m_t \gets \mathrm{cap}\!\left(\Delta\text{key-count}(x^m_t),\,\rho\right)$
           \COMMENT{per-range increment, capped by chunk size}
    \STATE Transfer fence: drain pending fast-tier P2P copies before any
           slow-tier transfer begins
    \FOR{$j = 1$ \TO $N$}
      \IF{$t < T_{\mathrm{warm}}$ \textbf{or} $t \bmod K_j = 0$}
        \COMMENT{transfer structure $s^j$ across the tier boundary}
        \STATE $s^{j,m}_t \gets \textsc{Update}^j\!\left(\mathbb{E}_m[s^{j}_{t-1}],\ \hat{g}^m_t\right)$
      \ELSE
        \STATE $s^{j,m}_t \gets \textsc{Update}^j\!\left(s^{j,m}_{t-1},\ \hat{g}^m_t\right)$
               \COMMENT{local step; cross-tier transfer elided}
      \ENDIF
    \ENDFOR
    \IF{$t < T_{\mathrm{warm}}$ \textbf{or} $t \bmod K_x = 0$}
      \COMMENT{re-balance and transfer the build partitioning}
      \STATE $x^m_{t+1} \gets \textsc{Part}\!\left(\textsc{HostOpt}(\mathbb{E}_m[x^m_t]),\ \hat{g}^m_t,\ \eta_t,\ \{s^{j,m}_t\}\right)$
    \ELSE
      \STATE $x^m_{t+1} \gets \textsc{Part}\!\left(x^m_t,\ \hat{g}^m_t,\ \eta_t,\ \{s^{j,m}_t\}\right)$
    \ENDIF
  \ENDFOR
\ENDFOR
\end{algorithmic}
\end{algorithm}

% --- Algorithm 1 prose description ---
% template: As shown in Algorithm 1, A syncs parameters x and states {s^j} at
%           state-specific intervals; setting N=2, s1=u, s2=v yields A-Adam
Motivated by the timescale separation above, we formalize AJB as a join pipeline
that composes a one-time topology-aware partitioning plan with a runtime
tier-routed transfer schedule. Algorithm~\ref{alg:ajb} specifies the complete
procedure. As shown there, AJB transfers the build partitioning $x \in
\mathbb{R}^d$ and the merge-path structures $\{s^j\}_{j=1}^{N}$ across the tier
boundary at structure-specific intervals $K_x, \{K_j\}_{j=1}^{N}$. Setting $N =
2$, $s^1_t = u_t$ (the boundary vector) and $s^2_t = v_t$ (the materialization
buffers), and instantiating the transfer rules with the running summaries of
Eq.~\eqref{eq:ema}, yields the concrete \emph{AJB--hybrid} instance evaluated in
this paper.

\paragraph{Planning phase.}
The plan is computed once, before any tuple is joined. The interconnect is
probed to classify every device pair as either a fast peer-to-peer tier or a
slow PCIe/host tier, and the build partitions are stratified into transfer tiers
by drift rate: dense, fast-changing key ranges form the $x$-tier and are
transferred at the base period; the merge-path boundary vector forms the
$u$-tier; and the materialization buffers, whose occupancy changes slowest, form
the $v$-tier. Device-local chunk sizes are set to occupy a fixed fraction of free
device memory and capped at an even share of the build relation, so that each
per-chunk key-count increment is bounded---the $|(g_t)_i| \leq \rho$ condition
that Eqs.~\eqref{eq:drift_u}--\eqref{eq:drift_v} rely upon.

\paragraph{Runtime phase.}
At each chunk group $t$, every device sorts (or radix-partitions) its local
chunk and immediately joins it against the co-resident probe range by binary
search---an intra-device step that is always executed because it is required for
the join's correctness. Upon completion, a \emph{transfer fence} drains any
pending fast-tier peer-to-peer copies before a slow-tier transfer may begin. This
fence is essential on heterogeneous fabrics: a device on the fast NVLink island
may otherwise enter a slow-tier transfer while a peer still has P2P copies in
flight on a shared link, violating the per-link ordering the runtime assumes and
risking deadlock. The slow-tier traffic is then gated by the tier-routed
schedule, with each structure's transfer independently triggered by its assigned
period:
\begin{itemize}[nosep,leftmargin=*]
\item \emph{$x$-tier} (build partitions): transferred every $K_x$ chunk groups.
\item \emph{$u$-tier} (merge-path boundary vector, whose drift follows the
  first-moment timescale): transferred every $K_u = 3K_x$ chunk groups.
\item \emph{$v$-tier} (materialization buffers, whose occupancy follows the
  second-moment timescale): transferred every $K_v = 6K_x$ chunk groups.
\end{itemize}
The ratios $K_u/K_x = 3$ and $K_v/K_x = 6$ approximately reflect the residency
half-lives of the corresponding summaries: the build partitioning is re-balanced
at the base rate, while the slowly drifting boundary and buffer structures
tolerate progressively longer local phases. During the warmup horizon
$T_{\mathrm{warm}}$, all tiers transfer at every chunk group ($K_x = K_u = K_v =
1$), recovering the uniform-cadence multi-GPU join for initial stability. Slow-tier
transfers are dispatched asynchronously on a dedicated copy stream, overlapping
cross-tier communication with the sort and join compute of the next chunk
group; full-width key accumulation is applied to the boundary counts before
transfer, avoiding the biased truncation drift that narrow counters would
introduce on long key ranges.

\paragraph{Composition guarantee.}
The two transfer families---intra-step fast-tier P2P copies and inter-step
slow-tier structure transfers---act along orthogonal axes and compose without
interference. The fast-tier copies ensure that each device's local join result
is \emph{complete} for its assigned key range; the decoupled slow-tier schedule
controls only \emph{how often} devices reconcile their partitionings with one
another. The correctness analysis of Section~\ref{sec:convergence} applies to
the join output produced after the fast-tier reconciliation, which is an
exact enumeration of the matches within each device's range; the AJB transfer
schedule then governs the cross-device merge cadence of these already-correct
partial results, and the guarantees carry through unchanged.

\paragraph{Relationship to the base join.}
The formulation applies generically to hash / sort-merge joins parameterized by
a partition rule $\textsc{Part}$ over $N$ merge-path structures $\{s^j_{-1}\}_{j=1}^{N}$,
each transferred by a rule $\textsc{Update}^j$. The chunk-size cap is
$[\mathrm{cap}(X,\rho)]_i = \mathrm{sgn}(X_i)\cdot\min\{|X_i|,\rho\}$. For a
provably balanced partitioning, $\textsc{HostOpt}$ reduces to the parallel
multiway-merge re-balancer on the CPU host; extension to a skew-adaptive host
re-balancer that biases boundaries toward dense ranges is deferred to the
appendix. Setting $N = 2$ as above and using the boundary / buffer update rules
yields AJB--hybrid; the construction also accommodates a radix-partitioning
build phase and a random-order enumeration variant of the merge path, the latter
drawing on randomized join enumeration~\citep{chen2023renum} to bound the cost of
re-deriving balanced boundaries after a device is added.

\paragraph{Illustrative example.}
To highlight the practical benefit of decoupled transfer, a toy
two-relation join under skewed, noisy key counts shows AJB and the
uniform-cadence P2P join both converging to the correct, balanced output, while
the pin-local heuristic---which fixes each partition to its first device---loses
load balance and the periodic-full-reset heuristic oscillates and overflows its
buffers. The corresponding figure and a non-uniform-key variant are
deferred to the evaluation and appendix.

%=============================================================================
\section{Correctness and Balance Guarantees}
\label{sec:convergence}
%=============================================================================
% Claude-2 (M026-M050): Section 3.
% Template anchor: \section{CI for A} in template_extraction_des_loc_reconstructed_EN.tex
% Mapping: CI=correctness/balance guarantees, CIV=partitioning problem,
%   CV/CVI/CVII=global/per-device/global imbalance objective, CVIII=device,
%   CIX=per-device key range, CX=uniform keys, CXVII=chunk size, CXVIII=chunk
%   limit eta_0, CXXII=psi cadence penalty, CXXIII=probabilistic transfer
%   equivalence, CXXV=finite period for the boundary vector.

This section establishes the theoretical foundation for AJB's decoupled
cross-tier transfer schedule. We analyze a single-structure variant of the join
under decoupled transfer cadences---tracking only the merge-path boundary
vector---and defer the full two-structure analysis (boundary vector and
materialization buffers) and the high-probability balance bounds to the
appendix. Throughout, ``correctness'' means that AJB enumerates \emph{exactly}
the set of matching tuple pairs of a single-device hash join, and
``balance'' refers to how evenly the build relation is partitioned across
devices.

A question specific to the composed pipeline is whether the
intra-step fast-tier peer-to-peer copies affect the analysis. We show that they
do not: the fast-tier reconciliation produces, on each device, an \emph{exact
local join} over that device's assigned key range, and the subsequent decoupled
slow-tier merge operates on these already-correct partial results. The analysis
therefore concerns only the cross-device partitioning cadence, and the output
cardinality is invariant to the choice of transfer periods.

\paragraph{Problem formulation.}
We cast the cross-device partitioning as the balance problem
\begin{equation}
\min_{x \in \mathbb{R}^d} f(x) := \frac{1}{M} \sum_{m=1}^{M} f_m(x),
\quad \text{with } f_m(x) = \mathbb{E}_{\xi \sim \mathcal{D}_m}[F_m(x; \xi)],
\label{eq:objective}
\end{equation}
where $M$ devices collaboratively minimize the global imbalance $f$ of a
candidate partitioning $x$ (a vector of per-device key-range assignments),
$f_m$ is the local imbalance contributed by device $m$, and $F_m(x;\xi)$ is the
per-chunk imbalance for a sampled key range $\xi$. In the general (skewed)
setting, each device $m$ observes keys drawn from a local distribution
$\mathcal{D}_m$ that may differ across devices; the uniform-key case
$\mathcal{D}_1 = \cdots = \mathcal{D}_M$ is recovered as a special instance.

When the join runs across a tier topology of degree $P$, each device processes
a contiguous chunk of the build relation and joins it against the co-resident
probe range by binary search. The fast-tier copies redistribute boundary
keys so that each device holds the exact build- and probe-side ranges bracketing
its partition. The resulting per-chunk imbalance increment $g_t^m =
\nabla F_m(x_t^m;\,\xi_t^m)$ is therefore computed with respect to the
\emph{same} balance objective $F_m$ as in the single-tier case---the fast-tier
redistribution introduces no additional bias or variance into the count
oracle.

\begin{figure}[t]
\centering
% [Figure 1 graphic: toy two-relation join; populated by a later Claude.]
\caption{(Left) Distance to a perfectly balanced partitioning and (right) the
2-D key-density contour for a toy skewed two-relation join over a hot key range,
with $M=256$ devices and IID Poisson count noise ($\sigma=1.5$).
AJB ($K_x=192$, $K_u=192$, $K_v=692$) and the uniform-cadence P2P join ($K=K_x$)
both converge to a balanced partition with identical output cardinality
(overlapping trajectories). The pin-local heuristic, which fixes each partition
to the device that first produced it (\protect\rule{6pt}{6pt}), fails to
balance. The periodic-full-reset heuristic (\protect\rule{6pt}{6pt}) stalls and
overflows its buffers due to repeated boundary oscillations. A non-uniform-key
variant appears in the appendix.}
\label{fig:toy}
\end{figure}

We impose the following standard regularity conditions on the balance problem.

\begin{assumption}[Lower bound and smoothness]
\label{asm:smooth}
The global imbalance $f\!:\mathbb{R}^d\!\to\!\mathbb{R}$ is bounded below by
some $f^*\!\in\!\mathbb{R}$ (the imbalance of the optimal partitioning), and
each local imbalance $f_m$ is $L$-smooth:
\begin{equation}
\|\nabla f_m(x) - \nabla f_m(y)\| \leq L\|x - y\|, \quad
\forall\, x,y \in \mathbb{R}^d.
\end{equation}
\end{assumption}

\begin{assumption}[Unbiased counts with bounded variance]
\label{asm:noise}
The per-chunk count increment $g^m$ measured by device $m$ is an unbiased
estimator of $\nabla f_m(x)$ with bounded variance:
\begin{equation}
\mathbb{E}[g^m] = \nabla f_m(x), \quad
\mathbb{E}[\|g^m_t - \nabla f_m(x)\|^2] \leq \sigma^2, \quad
\forall\, x \in \mathbb{R}^d.
\end{equation}
\end{assumption}

\begin{assumption}[Bounded skew]
\label{asm:heterogeneity}
For any $x \in \mathbb{R}^d$, the inter-device count heterogeneity satisfies
\begin{equation}
\frac{1}{M} \sum_{m=1}^{M} \|\nabla f_m(x)\|^2 \leq G^2 + B^2
\|\nabla f(x)\|^2.
\end{equation}
\end{assumption}

All three assumptions are standard in the analysis of partitioning and
load-balancing under
sampling~\citep{karimireddy2020scaffold,wang2021field,DeWitt1992practical}.
Setting $G^2 = 0$, $B^2 = 1$ recovers the uniform-key case.

\paragraph{Probabilistic transfer equivalence.}
To facilitate the analysis, we view each cross-tier transfer as a Bernoulli
trial: instead of re-balancing the build partitioning every $K_x$ chunk groups
($t \bmod K_x = 0$), we re-balance with probability $p_x = 1/K_x$ at each chunk
group. The two formulations are statistically equivalent in expectation and
permit a cleaner expression of the balance rate in terms of continuous transfer
probabilities.

\begin{theorem}[Balance rate of AJB under decoupled transfer]
\label{thm:sgdm}
Let Assumptions~\ref{asm:smooth},~\ref{asm:noise}, and~\ref{asm:heterogeneity}
hold. Choose the device-local chunk size $\eta = \min(\eta_0,\, T^{-1/2})$ with
\begin{equation}
\eta_0 \stackrel{\mathrm{def}}{=}
\frac{1}{4L} \min\!\left(1 - \beta,\;
\frac{1}{6\sqrt{\psi\,\max(1,\, B^2 - 1)}}\right),
\quad \text{where}\;\;
\psi \stackrel{\mathrm{def}}{=}
\frac{4(1-p_x)}{p_x^2}
\cdot \frac{(1-\beta)(1-p_u)}{1-(1-p_u)\beta}\,.
\label{eq:step_size}
\end{equation}
Then the averaged partitioning $\bar{x}_t = \mathbb{E}_m[x^m_t]$ produced by AJB
approaches a balanced assignment at rate
\begin{equation}
\frac{1}{T}\sum_{t=0}^{T-1}
\mathbb{E}\|\nabla f(\bar{x}_t)\|^2
\;\leq\; \frac{4}{\sqrt{T}}
\left(f(x_0) - f^* + \frac{L\sigma^2}{2M}\right)
+ \mathcal{O}\!\left(\frac{1+\psi}{T}\right).
\label{eq:rate}
\end{equation}
\end{theorem}

We now discuss the implications of this result and its interaction with the
composed join pipeline.

\paragraph{Asymptotic optimality.}
The leading $\mathcal{O}(1/\sqrt{T})$ term in~\eqref{eq:rate} matches the
minimax lower bound for stochastic balancing under count
noise~\citep{arjevani2023lower} and is \emph{independent of the number of local
chunks processed between transfers}. The transfer probabilities $p_x$, $p_u$ and
the drift parameter $\beta$ enter only the higher-order $\mathcal{O}(1/T)$ term
through the cadence-penalty factor $\psi$, and therefore exert a limited impact
on the asymptotic balance rate. The $1/M$ factor on the variance term confirms a
\emph{linear speedup}: aggregating counts across more devices proportionally
sharpens the balance estimate.

\paragraph{Asymmetric sensitivity to build-partition vs.\ boundary transfer.}
The structure of $\psi$ reveals a fundamental asymmetry between transferring the
build partitioning and transferring the merge-path boundaries. Build-partition
transfer exerts a \emph{quadratic} influence: $\psi = \mathcal{O}(1/p_x^2)$, so
that letting $p_x \to 0$ (never re-balancing) renders the imbalance term
unbounded and breaks the balance guarantee. Boundary transfer, by contrast,
enters \emph{linearly}: setting $p_u = 0$ (never transferring the boundary
vector) does not destroy the asymptotic rate. However, the chunk-size
restriction~\eqref{eq:step_size} depends on $\sqrt{\psi}$, and $\psi$ is
maximized as $p_u \to 0$ (since $(1-p_u)/(1-(1-p_u)\beta) \to 1/(1-\beta)$),
imposing the most severe constraint on the admissible chunk size $\eta_0$.

This analysis yields the central design principle underlying AJB's tiered
transfer schedule:
\begin{quote}
\emph{The build partitioning must be transferred frequently ($p_x$ large,
equivalently $K_x$ small). The merge-path structures may be transferred less
often---but increasing the boundary-transfer frequency $p_u$ relaxes the
chunk-size restriction, enabling larger device-local chunks, greater copy/compute
overlap, and faster practical balancing.}
\end{quote}

\paragraph{Interaction with tier-routed transfer.}
In AJB, the theoretical distinction between build-partition and boundary
transfer frequencies is operationalized through the tier-routed scheduler of
Section~\ref{sec:ajb_algorithm}. The $x$-tier (build partitions) is transferred
at the base period $K_x$, reflecting the quadratic sensitivity
$\psi \propto 1/p_x^2$. The $u$-tier (merge-path boundary vector) and $v$-tier
(materialization buffers) are transferred at $K_u = 3K_x$ and $K_v = 6K_x$
respectively, reflecting the empirical finding that the boundary and buffer
structures tolerate longer local phases in proportion to their residency
half-lives (Section~\ref{sec:timescales}). The warmup phase
$t < T_{\mathrm{warm}}$ forces $K_x = K_u = K_v = 1$, recovering the
uniform-cadence multi-GPU join until the partitioning enters a stable basin.

\paragraph{Composition with fast-tier reconciliation.}
A natural concern is whether the intra-step fast-tier copies introduce
additional terms into the balance bound. They do not, because the fast-tier
reconciliation is a \emph{deterministic, lossless} redistribution of boundary
keys: each device joins its exact assigned build- and probe-side ranges,
producing precisely the matches a single-device join would enumerate (up to the
order in which materialized rows are appended, which does not affect
cardinality). The per-chunk count increment $g_t^m$ entering the decoupled
transfer loop of Algorithm~\ref{alg:ajb} is therefore an unbiased estimator of
$\nabla f_m(x_t^m)$ with variance bounded by $\sigma^2$, and
Assumptions~\ref{asm:noise}--\ref{asm:heterogeneity} hold with the same
constants as in the single-tier case.

\paragraph{Extension to buffers and high-probability bounds.}
Our analysis extends to the full two-structure join with decoupled boundary- and
buffer-transfer periods under mildly skewed key distributions (appendix). The
high-probability analysis establishes that---unlike the in-expectation bound
above, which tolerates $p_u = 0$---the merge-path boundary vector must be
transferred with \emph{finite} period whenever the key distribution is
non-uniform ($\beta_2 < 1.0$), i.e., $K_v < \infty$. This is consistent with the
residency half-life intuition: a buffer-occupancy summary with $\beta_2 = 0.999$
has $\tau_{0.5} \approx 693$ chunks, which is large but finite; indefinitely
postponing its transfer lets the per-device occupancy estimates drift unboundedly,
risking buffer overflow and violating the high-probability tail condition that
guarantees a non-overflowing merge path.

%=============================================================================
\section{Experimental Design}
\label{sec:experimental_design}
%=============================================================================
% Claude-3 (M051-M075): Section 4 + Section 5 lead-in.
% Template anchor: \section{Experimental Design} in
%   template_extraction_des_loc_reconstructed_EN.tex
% Mapping: RQ1-RQ6 re-grounded on relation-resident structures; CLXX/CLXXIV=
%   13B-tuple scale; CLX=chunk cap rho; CLXIII=K_x=256; CLXVII=A6000x2 + H100 NVL
%   + CPU host; CLIV=2x; CXLV=skew-adaptive host re-balancer; CXLVII=random-order
%   merge-path enumeration. Metrics: end-to-end join time, slow-tier byte volume,
%   relative structure drift, output cardinality (correctness).

Our experimental evaluation addresses the following research questions:
\begin{description}[nosep,leftmargin=*]
\item[RQ1] Do the theoretical drift rates (Section~\ref{sec:timescales})
  predict the empirical evolution of the relation-resident structures across
  device-local chunks?
\item[RQ2] How does the cross-tier transfer period of each structure class---
  build partitions, merge-path boundaries, materialization buffers---impact
  end-to-end join time?
\item[RQ3] To what extent can AJB reduce slow-tier byte volume relative to the
  uniform-cadence P2P join in practical join configurations?
\item[RQ4] How does AJB scale with increasing relation size and toward
  out-of-core, billion-tuple joins?
\item[RQ5] How does the join perform when a skew-adaptive host re-balancer
  replaces plain multiway-merge averaging?
\item[RQ6] Is the decoupled transfer schedule compatible with non-sort-merge
  build phases, including a random-order merge-path enumeration variant?
\end{description}

%-----------------------------------------------------------------------------
\subsection{Experimental Setup}
\label{sec:experimental_setup}
%-----------------------------------------------------------------------------

\paragraph{Relations and data.}
We join a build relation $R$ against a probe relation $S$ of $8$-byte
key/$8$-byte payload tuples, generated under controlled key distributions. We
distinguish the per-device chunk size $B_w$---capped at a fixed fraction of free
device memory---from the global input $B = \sum_{m=0}^{M-1} B_w$ resident across
the $M$ devices. Inputs range from $1$\,B to $13$\,B tuples and are sampled under
both a uniform key distribution and a Zipfian distribution with skew parameter
$\theta$ controlling the fraction of tuples falling in the hottest key range,
following the generator conventions of the multi-GPU join
benchmark~\citep{maltry2025multigpu}. Match selectivity is fixed by a target
output cardinality so that the materialization volume is comparable across
distributions. Foreign-key and many-to-many join variants---in which a single
hot key fans out to a large materialized range---are reported in the appendix.

\paragraph{Join configurations.}
We employ a sort-merge build phase backed by parallel CPU
sorting~\citep{kim2009sortvshash} and a hybrid-radix build phase that partitions
on the device, the latter modifying the partitioning rule to remain balanced
under any key skew. For the $1$\,B-tuple workload we grid-search the chunk-size
cap, the base transfer period $K_x$, and the device tier assignment; larger
workloads reuse the resulting configuration. Chunk sizes follow a
warmup-stable-decay schedule that synchronizes all tiers during the warmup
horizon before relaxing to the decoupled periods. We favor the hybrid-radix
build with a slowly drifting buffer-occupancy summary ($\beta_2 = 0.9999$) in
high-skew regimes, where eager re-partitioning otherwise thrashes the slow tier.
We also ablate the host re-balancer, comparing plain multiway-merge averaging
against a skew-adaptive re-balancer that biases boundaries toward dense ranges,
on the $1$\,B-tuple workload.

\paragraph{Baselines.}
We compare AJB with: (i)~eager per-level repartitioning, which reshuffles every
structure across the device boundary at every partitioning level (upper bound on
balance quality, worst case on slow-tier volume); (ii)~the uniform-cadence
P2P-enabled multi-GPU join with a single transfer period
$K$~\citep{maltry2025multigpu}; (iii)~the pin-local heuristic, which fixes every
structure to the device that first produced it; and (iv)~the periodic-full-reset
heuristic, which recomputes all boundary structures from scratch after each
phase. The pin-local heuristic corresponds to AJB with $K_u, K_v = \infty$ and
constitutes an upper bound on slow-tier efficiency; eager repartitioning
constitutes an upper bound on partition-balance quality.

\paragraph{Metrics.}
We evaluate by (i)~end-to-end join time, decomposed into its sort, merge, and
join phases, and (ii)~per-device asymptotic slow-tier byte volume, assuming a
bandwidth-optimal peer-to-peer ring for the fast tier and scaling linearly with
the transferred structure size. We additionally report (iii)~output cardinality,
which must match the single-device join exactly and serves as our correctness
check across all transfer cadences. To fairly compare structure evolution across
drift rates, we measure the \emph{relative} rate of change
$\|s_{t+K} - s_t\|_2 / \|s_t\|_2$. We report both stepwise (per chunk group) and
wall-clock convergence to a balanced partitioning. Our hardware consists of a
single heterogeneous server combining two RTX A6000s (NVLink-bridged fast tier)
and one H100 NVL with a multi-core CPU host; the cross-tier host link runs at
$100$\,Gb/s with practical throughput of $60$--$70$\,Gb/s. An analysis of join
time versus bandwidth appears in the appendix.

%=============================================================================
\section{Evaluation}
\label{sec:evaluation}
%=============================================================================

Our results demonstrate that the relation-resident structures evolve at
empirically distinct rates governed by their drift constants
(Section~\ref{sec:rq1}), forming a clear transfer hierarchy
(Section~\ref{sec:rq2}). AJB reduces slow-tier byte volume by $2\times$ relative
to the uniform-cadence P2P join (Section~\ref{sec:rq3}), converges robustly to a
balanced partition under device addition, and scales effectively toward
billion-tuple joins (Section~\ref{sec:rq4}). The decoupled schedule composes with
both a skew-adaptive host re-balancer (Section~\ref{sec:rq5}) and a random-order
merge-path enumeration variant of the build phase (Section~\ref{sec:rq6}).

%-----------------------------------------------------------------------------
\subsection{Higher-Drift Structures Exhibit Slower Empirical Rates of Change (RQ1)}
\label{sec:rq1}
% Claude-4 (M076-M100): RQ1. Template anchor: \subsection{RQ1: Empirical Rate
% of Change}. Mapping: momenta->boundary vector / buffer occupancy; CLX=rho=1;
% CLXI=half-life ratio. Figure 2 = relative drift figure.
%-----------------------------------------------------------------------------

Figure~\ref{fig:rates_of_change} confirms that the relative rates of change of
the two merge-path structures---the boundary-count vector and the
buffer-occupancy summary---scale with their respective drift rates under the
bounded per-chunk count increment ($\rho = 1$). Consistent with the residency
half-life analysis of Section~\ref{sec:timescales}, the buffer-occupancy summary
evolves substantially slower than the boundary vector whenever $\beta_2 \gg
\beta_1$. The occupancy estimate remains slower even when $\beta_2 \approx
\beta_1$, because the squared-count variance underlying the buffer summary
accumulates on an inherently slower timescale than the raw boundary counts that
drive the boundary vector.

\paragraph{Takeaway.}
When $\beta_1 \ll \beta_2$, the buffer-occupancy summary's empirical rate of
change is slower than the boundary vector's by a factor proportional to the
residency half-life ratio $\tau_{0.5}(\beta_2)/\tau_{0.5}(\beta_1) =
\ln(\beta_1)/\ln(\beta_2)$, validating the motivation for assigning a longer
cross-tier transfer period to the slower-drifting structure.

\begin{figure}[t]
\centering
% [Figure 2 graphic; populated by a later Claude from the relative-drift data.]
\caption{Relative rates of change for the boundary-count vector (left) and the
buffer-occupancy summary (right) across transfer rounds for the hybrid-radix
build ($K = 64$). At $\beta_2 = 0.9999$, increasing $\beta_1 \geq 0.99$ greatly
slows the boundary vector's rate of change. The buffer-occupancy summary evolves
${\sim}100\times$ slower, consistent with its residency half-life being two
orders of magnitude larger.}
\label{fig:rates_of_change}
\end{figure}

%-----------------------------------------------------------------------------
\subsection{Build Partitions Require Frequent Transfer; Boundary Transfer Scales
with Half-Life (RQ2)}
\label{sec:rq2}
% Template anchor: \subsection{RQ2: Independent R}. CLXIII=K_b=16/256.
% Figure 3 = independent-period sweep.
%-----------------------------------------------------------------------------

Figure~\ref{fig:sync_freq} evaluates the effect of independently varying the
transfer periods $K_x$, $K_u$, and $K_v$. We consider two base periods ($K_b =
16$ and $K_b = 256$), chosen to bracket the fastest structure's residency
half-life ($\tau_{0.5}(0.95) \approx 13.5$). Frequent transfer of the build
partitioning ($K_x$) is critical for balance quality and end-to-end time,
consistent with the quadratic sensitivity $\psi \propto 1/p_x^2$ established in
Theorem~\ref{thm:sgdm}. By contrast, the transfer frequency of the merge-path
structures ($K_u$, $K_v$) materially affects join time only when the assigned
period is commensurate with the corresponding structure's half-life; outside
this resonance regime, boundary and buffer transfer frequency primarily
influences slow-tier volume rather than balance quality.

\paragraph{Takeaway.}
Build-partition transfer frequency dominates balance quality and join time, in
agreement with the leading-term analysis of Section~\ref{sec:convergence}.
Merge-path transfer periods matter empirically only when chosen near their
residency half-lives, validating the timescale separation principle of
Sections~\ref{sec:timescales} and~\ref{sec:convergence}.

\begin{figure}[t]
\centering
% [Figure 3 graphic; populated by a later Claude.]
\caption{End-to-end join time for AJB (hybrid-radix, $\beta_1 = 0.95$,
$\beta_2 = 0.9999$), varying transfer periods independently while holding the
others at $K_b$. (a)~Build-partition transfer is critical, with clear
degradation at higher $K_x$. (b)~Buffer transfer minimally affects join time due
to its large residency half-life ($\tau_{0.5} \approx 6931$). (c,d)~Boundary
transfer improves balance when the period matches its half-life; contrast the
two base regimes. Results for the sort-merge build appear in the appendix.}
\label{fig:sync_freq}
\end{figure}

%-----------------------------------------------------------------------------
\subsection{AJB Achieves $2\times$ Cross-Tier Volume Reduction over the
Uniform-Cadence P2P Join (RQ3)}
\label{sec:rq3}
% Template anchor: \subsection{RQ3: N Reduction}. CLIV=2x, CLXV=1.37x, CLXVI=
% 1.1x, CLXVII=hardware. Figure 4 = comms reduction + device addition.
%-----------------------------------------------------------------------------

As shown in Figure~\ref{fig:comms}, AJB halves the slow-tier byte volume
relative to the uniform-cadence P2P join~\citep{maltry2025multigpu} at matching
output cardinality and balance, by transferring the merge-path structures less
frequently ($K_u = 3K_x$, $K_v = 6K_x$) and exploiting the buffer summary's
lower sensitivity to transfer frequency (Figure~\ref{fig:sync_freq}). At $K_x =
16$ on the A6000$\times$2 $+$ H100 server, this yields a $1.37\times$ end-to-end
speedup over eager per-level repartitioning and $1.1\times$ over the
uniform-cadence join. At $K_x = 256$, the speedup over eager repartitioning
increases to $1.47\times$, and both AJB and the uniform-cadence join approach the
fast-tier-only throughput ceiling at which the slow link is no longer the
bottleneck.

Figure~\ref{fig:comms}(c,d) further demonstrates that AJB cleanly admits a newly
attached device via its periodic structure transfer, re-deriving a balanced
partitioning and outperforming the ad-hoc alternative of reusing a stale,
device-local boundary vector---whose insufficiency under skew is detailed in the
appendix.

\paragraph{Takeaway.}
AJB achieves a $2\times$ reduction in slow-tier byte volume over the
uniform-cadence join by leveraging two complementary insights from the balance
analysis: (i)~merge-path transfer has a weaker asymptotic impact than
build-partition transfer ($\psi$ depends linearly on $p_u$ but quadratically on
$1/p_x$), and (ii)~structures with larger residency half-lives tolerate longer
local phases without balance degradation. By eventually transferring all
structures at finite periods, AJB matches the device-fault-tolerance of the
uniform-cadence join---with effective period $K = \max(K_x, K_u, K_v)$---when
incorporating a new device or recovering from a stalled tier.

\begin{figure}[t]
\centering
% [Figure 4 graphic; populated by a later Claude.]
\caption{AJB ($K_u = 3K_x$, $K_v = 6K_x$) reduces slow-tier byte volume by
$2\times$ over the uniform-cadence P2P join while matching its output cardinality
and balance, and that of the heuristic baselines, in both the high-frequency (a)
and low-frequency (b) regimes. Panels (c,d) demonstrate robustness: attaching an
additional device at chunk group $1536$, AJB and the uniform-cadence join
maintain a stable, balanced partitioning, outperforming the pin-local and
periodic-reset heuristics and stale-boundary reuse.}
\label{fig:comms}
\end{figure}

%-----------------------------------------------------------------------------
\subsection{AJB Scales Toward Billion-Tuple, Out-of-Core Joins (RQ4)}
\label{sec:rq4}
% Template anchor: \subsection{RQ4: Large-Scale Training}. CLXX/CLXXIV=13B,
% CLXVIII=256, CLXXI=Table 1 (quality), CLXXIII=2.2x. Figure 5 + Table 1/2.
%-----------------------------------------------------------------------------

Figure~\ref{fig:billion} demonstrates that AJB reliably scales toward
billion-tuple, out-of-core joins under highly infrequent transfer ($K_x = 256$).
Evaluating on out-of-core join quality across input distributions
(Table~\ref{tab:icl}), AJB matches the output cardinality and balance of all
baselines while moving substantially fewer bytes across the slow tier than both
the uniform-cadence join and eager repartitioning. The pin-local
heuristic~\citep{maltry2025multigpu} suffers from load imbalance under
skew---manifesting as buffer-occupancy growth on the hot device
(Figure~\ref{fig:billion}b)---which inflates its tail join time and underscores
the advantage of AJB's principled transfer schedule. We elaborate on the
conditions under which decoupled transfer improves stability under skew in the
appendix.

The reduced slow-tier volume translates to a ${\approx}\,2.2\times$ end-to-end
join-time speedup over eager per-level repartitioning (Figure~\ref{fig:billion}).
As shown in Table~\ref{tab:wallclock}, these savings scale with input size and
transfer period. At the $13$\,B-tuple scale with $K_x = 16$, AJB saves more than
$13$ time units relative to the uniform-cadence join and over $73$ relative to
eager repartitioning (in the normalized units of Table~\ref{tab:wallclock}). The
advantage over eager repartitioning widens further at $K_x = 256$, where the
cross-tier-efficient methods approach peak device utilization. Equivalent
throughput figures appear in the appendix.

\begin{figure}[t]
\centering
% [Figure 5 graphic; populated by a later Claude.]
\caption{AJB matches the uniform-cadence join's time-to-balance (left) for
billion-tuple joins at half the slow-tier cost ($K_x = 256$, $K_u = 3K_x$,
$K_v = 6K_x$), representing a $\mathbf{170\times}$ reduction over eager per-level
repartitioning. Both AJB and the uniform-cadence join produce identical output
cardinality at this scale. The pin-local heuristic reaches comparable
time-to-balance on uniform keys (left) but exhibits buffer-occupancy growth on
the hot device (right) and partition drift under skew, raising concerns for
adversarial workloads.}
\label{fig:billion}
\end{figure}

\begin{table}[t]
\centering
\caption{Billion-tuple out-of-core join quality across input distributions
(higher is better; balance score, normalized so that a perfectly even partition
$= 100$). AJB matches or surpasses the partition balance achieved under the
uniform-cadence join and the pin-local heuristic, approaching eager
repartitioning. The pin-local heuristic underperforms under the skewed
(Zipfian, foreign-key) distributions, consistent with the buffer-occupancy
growth of Fig.~\ref{fig:billion}b, indicating that stale device-local boundaries
have unbalanced the partitioning.}
\label{tab:icl}
\begin{tabular}{lcccc|c}
\toprule
& Uniform & Zipfian & Foreign-key & Many-to-many & Avg \\
\midrule
AJB                & 31.8 & 59.0 & 70.7 & 44.9 & 51.6 \\
Uniform-cadence    & 31.9 & 59.0 & 70.6 & 45.8 & 51.8 \\
Pin-local          & 30.1 & 58.0 & 70.0 & 44.8 & 50.7 \\
\midrule
Eager repartition  & 33.8 & 62.5 & 71.1 & 47.8 & 53.8 \\
\bottomrule
\end{tabular}
\end{table}

\begin{table}[t]
\centering
\caption{End-to-end join time (normalized units) for $1$--$13$\,B-tuple inputs at
high ($K_x = 16$) and low ($K_x = 256$) transfer frequencies. At high frequency,
AJB outperforms the uniform-cadence join by over $13$ units on the $13$\,B input
and is within $3\%$ of the pin-local heuristic. At low frequency, it reduces the
$13$\,B join time by $>93$ units versus eager repartitioning.}
\label{tab:wallclock}
\begin{tabular}{l|cc|cc|cc}
\toprule
& \multicolumn{2}{c|}{1B tuples} & \multicolumn{2}{c|}{7B tuples}
& \multicolumn{2}{c}{13B tuples} \\
$K_x$ & 16 & 256 & 16 & 256 & 16 & 256 \\
\midrule
Eager repartition & $1.41 \pm 0.008$ & $1.41 \pm 0.008$
& $38.74 \pm 0.161$ & $38.74 \pm 0.161$
& $175.50 \pm 0.478$ & $175.50 \pm 0.478$ \\
\midrule
Pin-local & $0.80 \pm 0.007$ & $0.63 \pm 0.006$
& $28.52 \pm 0.095$ & $24.01 \pm 0.088$
& $100.21 \pm 0.544$ & $82.46 \pm 0.513$ \\
Uniform-cadence & $0.96 \pm 0.006$ & $0.64 \pm 0.006$
& $31.46 \pm 0.090$ & $24.18 \pm 0.087$
& $116.10 \pm 0.484$ & $83.34 \pm 0.509$ \\
AJB & $0.81 \pm 0.006$ & $0.63 \pm 0.006$
& $28.80 \pm 0.094$ & $24.06 \pm 0.088$
& $102.03 \pm 0.537$ & $82.68 \pm 0.512$ \\
\bottomrule
\end{tabular}
\end{table}

\paragraph{Takeaway.}
AJB enables efficient billion-tuple joins, particularly on out-of-core inputs
where the slow tier dominates, with partition balance and output cardinality
competitive with eager repartitioning. We recommend setting $K_x$ to balance
throughput against the available cross-tier bandwidth, then choosing $K_u$,
$K_v$ as fixed multiples of $K_x$ (e.g., $3\times$, $6\times$) or calibrating
them against the residency half-lives of their respective $\beta$ values
(see appendix).

%=============================================================================
% Section 5.5-5.6 (RQ5-RQ6) and Sections 6-7 are left for subsequent Claudes
% per the development plan. Symbol mapping is fixed in ajb_mapping.txt.
%=============================================================================

\end{document}
